%%%%%%%%%%  scan idx 43  |  folio 30  %%%%%%%%%%
\begin{equation}\tag{1.19}
\eta'(x^i)=(-1)^{i-1}(i-1)!\,x^i\qquad (x^i\in A^i)
\end{equation}
est un \emph{homomorphisme d'anneaux}.

On introduit dans $\widetilde A$ des opérations $\lambda^i$ vérifiant les axiomes des $\lambda$-anneaux (\S 1, formules (1.2)), bien déterminées par la condition que les composantes de $\lambda^i[0,x]$ se déduisent des composantes $x^i$ de $x$ par des polynômes \emph{universels} à coefficients entiers
\begin{equation}\tag{1.20}
\lambda^i[0,x]=[0,\lambda^i x]\qquad (\lambda^i x)^{(n)}=Q_{i,n}(x^1,\ldots,x^n)
\end{equation}
($Q_{i,n}$ est isobare de poids $n$), et que l'on ait
\begin{equation}\tag{1.21}
\lambda^i[1,1+x^1]=0\qquad\text{pour }i>1.
\end{equation}
Les polynômes $Q_{i,n}$ s'évaluent par un calcul de fonctions symétriques élémentaires, qui est le calcul des classes de Chern des puissances extérieures d'un fibré vectoriel, à l'aide des classes de Chern et du rang de ce dernier. On constate alors que $\widetilde A$ est un $\lambda$-anneau \emph{spécial} augmenté et que les $\widetilde A^i$ sont stables par les opérations $\lambda^j$.

Le $\lambda$-anneau augmenté $\widetilde A$ peut être considéré comme un foncteur en $A$~; par suite l'automorphisme principal
\[
x=\sum_{i\geq0}x^i\longrightarrow \check x=\sum_{i\geq0}(-1)^i x^i
\]
de $A$ définit un automorphisme de $\widetilde A$, noté par le même symbole $\check{\ }$.

\textbf{Proposition 1.2.} \emph{Pour $x\in\widetilde A$ de la forme $[n,1+x^1+\cdots+x^n]$} (on a donc $x^i=0$ pour $i>n$), \emph{on a $\lambda^i(x)=0$ pour $i>n$, on a $\lambda^n(x)=[1,1+x^1]$ et enfin}
\begin{equation}\tag{1.22}
\sum_{i=0}^n(-1)^i\lambda^i(x)=[0,1-(n-1)!x^n+\cdots]
\end{equation}
De cette proposition et du fait que les $\widetilde A^i$ sont stables par les opérations $\lambda^j$, on déduit que la formule (1.22) est valable chaque fois que $\varepsilon(x)=n$, que l'on ait ou non $x^i=0$ pour $i>n$.

